%% title: Next human ratings by judge-ensemble uncertainty, and the power curve for the held-out RPS gain (11-27-2026)
%% result-date: 11-27-2026
%% result-time: 11:23
%% experiment: 30_rank_next_human_ratings_by_judge_ensemble_uncertainty_and_power_curve_for_the_held_out_rps_gain_start_27Nov2026
%% headline: Top 25 new artifacts are 21/61 fives against 66/70 in the bottom 25. Extra ratings on the same 20 tasks leave the RPS upper bound at +0.002. A planning extrapolation crosses at 23 tasks.
% Source: experiments/30_rank_next_human_ratings_by_judge_ensemble_uncertainty_and_power_curve_for_the_held_out_rps_gain_start_27Nov2026/results.md

\subsection*{Bottom line}
The next human ratings should go to artifacts the six logged judges do not already call a five.
On the 100 new artifacts the top 25 are 21/61 fives (34.4\%) and the bottom 25 are 66/70 fives (94.3\%), against 185/250 fives (74.0\%) in the new set.
Replaying the experiment 27 head on the 445 ratings (195 old + 250 new) gives test ranked probability score (RPS) 0.0737 versus 0.1088 for the set-aware constant, paired $-0.0351$ $[-0.0764,+0.0020]$ (10,000 task draws, seed 20261127).
Cross-entropy (CE) is 0.943 versus 1.159, paired $-0.2153$ $[-0.4593,+0.0119]$.
Total variation distance (TVD) is 0.359 versus 0.479, paired $-0.1206$ $[-0.2038,-0.0506]$.
Spearman 0.876 and Pearson 0.907 are secondary.
Bias is $+0.213$ (judge mean 3.913 versus human mean 3.700).
Simulated extra ratings on these same 20 tasks, through $+2{,}240$ draws from the current histograms, leave the median RPS upper bound within 0.0008 of $+0.0020$.
It does not go negative.
An extrapolation that holds the point estimate fixed and scales only the upper distance by $\sqrt{20/m}$ is negative at 23 tasks.
That extrapolation is planning, not a certification.

\subsection*{Held-out test, 20 tasks, 70 artifacts (lower is better for RPS, CE, TVD)}
\begin{center}\scriptsize\setlength{\tabcolsep}{3pt}
\FitWidth{\begin{tabular}{@{}l l l l l l@{}}
\toprule
Forecast & RPS [95\%] & CE [95\%] & TVD [95\%] & Spearman & Pearson \\ \midrule
F8 refit & 0.0737 [0.0462, 0.1062] & 0.943 [0.665, 1.265] & 0.359 [0.292, 0.439] & 0.876 & 0.907 \\
Set-aware constant & 0.1088 [0.0725, 0.1503] & 1.159 [0.866, 1.468] & 0.479 [0.402, 0.563] &  &  \\
\bottomrule
\end{tabular}}
\end{center}

The experiment 27 replay on the frozen 247-new-rating table matches its published paired RPS at $-0.034895$ $[-0.075959,+0.002443]$ against the recorded target $-0.034895$ $[-0.075959,+0.002444]$.

\subsection*{Caveats}
\begin{itemize}
  \item Every power number is a retrospective subsample or a simulation from the current histograms. It is not a new human round and not a certification.
  \item The task bootstrap keeps the histograms fixed. More ratings drawn from those histograms do not shrink it.
  \item The 23-task figure assumes new tasks would have the same paired RPS. No logged six-member ensemble exists here for VeriBench tasks outside these 175 artifacts.
  \item No differential testing and no Lean check. The rank does not use rating values or the train, validation, or test split.
\end{itemize}

\subsection*{Appendix: methods, harness and model}
No new model call.
Logged scores are the H2 Astra and H2 Opus calls, the H3 Astra call, and the H4 Astra, Argon and Opus calls already stored in experiments 27, 28 and 29.
The head is the experiment 27 F8 cumulative logit, cross-entropy, penalty 0, mixture 0, refit on train and validation with NumPy 1.26.4 and SciPy 1.13.1.
Receipt: the experiment \path{results.md}.
Google Doc: \url{https://docs.google.com/document/d/1Nvwklzb89u223bSt5pU92F9lN7yf6ZvUJYOsYK1DpIg/edit}.
Priority list: \path{results/next_ratings_priority_top25_11-27-2026.csv} in that experiment.
